%% Énergie cumulée e(t) sur 24 h : deux rampes de pente 2 kW séparées et
%% suivies de paliers (pente nulle). Points imposés :
%% (0;0) -> (6;12) -> (12;12) -> (18;24) -> (24;24).
\documentclass[tikz,border=4pt]{standalone}
\usepackage{liaisons}
\begin{document}
\begin{tikzpicture}[x=0.25cm, y=0.145cm,
  axe/.style={-{Stealth[length=5pt]}, line width=0.6pt},
  courbe/.style={solideD, line width=1.4pt, line join=round, line cap=round},
  grad/.style={font=\scriptsize},
  repere/.style={line width=0.5pt},
  aide/.style={black!25, line width=0.4pt, dotted}]

%% ---- axes ---------------------------------------------------------------
\draw[axe] (-1.5,0) -- (27,0) node[anchor=south east, inner sep=4pt, font=\small] {$t$ (h)};
\draw[axe] (0,-2) -- (0,29);
\node[anchor=south west, inner sep=2pt, font=\small] at (0.5,27) {$e(t)$ (kWh)};

%% ---- graduations --------------------------------------------------------
\foreach \e in {12,24} {
  \draw[repere] (0,\e) -- (-0.7,\e) node[grad, anchor=east, inner sep=2pt] {\e};
  \draw[aide] (0,\e) -- (24,\e); }
\node[grad, anchor=north east, inner sep=2pt] at (-0.2,0) {0};
\foreach \t in {6,12,18,24} {
  \draw[repere] (\t,0) -- (\t,-1) node[grad, anchor=north, inner sep=2pt] {\t}; }

%% ---- courbe -------------------------------------------------------------
\draw[courbe] (0,0) -- (6,12) -- (12,12) -- (18,24) -- (24,24);

%% ---- annotations de pente ------------------------------------------------
\node[font=\scriptsize, text=solideB, anchor=south, inner sep=3pt, rotate=49] at (3,6)
  {pente = 2 kW};
\node[font=\scriptsize, text=solideB, anchor=south, inner sep=4pt] at (9,12)
  {pente nulle : 0 kW};
\end{tikzpicture}
\end{document}
